\section{The Inferential Cause and the Opening for Freedom}
\label{p13:sec:freedom}

\noindent
The jurisprudential chapter ended with a distinction it had traced through eight positions but not yet examined: between the mechanical cause, in which the anchor of validity sits at the beginning of the law's movement and at its abandoned limits, and the inferential cause, toward which that movement has been travelling. The distinction was drawn, in the diagnosis, as a contrast between two ways a result might be produced\,---\,the compulsion of a body and the grounding of a mind's expectation. This section asks what the inferential cause \emph{is}, and answers that it is not a weaker species of the mechanical but a different kind altogether, and that the difference is the very opening in which freedom, and with it the moral, becomes possible. The section is short, but it is the load-bearing arch beneath everything the mechanism, the justice criterion, and the praxis will build: it establishes that to generate trust is to address a free being as free, and that the promise which carries no force is, for exactly that reason, not the defective case but the proper one.

\subsection{Two kinds of cause, not two degrees of one}
\label{p13:sub:free-twokinds}

It is tempting to hear the contrast between mechanical and inferential cause as a contrast of strength\,---\,as though the mechanical cause were the reliable kind, which compels its effect outright, and the inferential cause a weaker kind, which merely makes its effect more probable. On that hearing, force would be the gold standard of causation and trust a degraded approximation to it, and the whole movement the jurisprudence traced would be a movement from a strong ground to a weak one, a story of decline. This is the hearing the paper must refuse, and the refusal is the hinge of the section.

The mechanical cause and the inferential cause are not two degrees of one kind of causation; they are two kinds. The mechanical cause operates upon its effect from outside and leaves the effect no part to play: the struck ball does not consent to move, contributes nothing to its own motion, could not have done otherwise given the impact. Its paradigm is the transfer of momentum, and its mark is that the affected thing is passive throughout\,---\,the cause does all the work, and the effect is the mere registration of it. The inferential cause operates altogether differently. It does not push; it gives a ground. It supplies a mind with something on the basis of which that mind may form an expectation\,---\,an expectation the mind itself forms, by its own activity, and could, being what it is, have formed otherwise or withheld. The ground does not compel the expectation as the impact compels the motion; it furnishes a reason, and a reason is precisely what a free mind takes up rather than undergoes. Where the mechanical effect is the passive registration of a force, the inferential effect is the active production of an expectation by the one in whom it arises.

This is why the inferential cause is not a weaker mechanical cause but a stronger\,---\,or rather a different and higher\,---\,thing. It is the only kind of causation that can have a free being for its proper subject, because it is the only kind that addresses that being \emph{as} free: as one who forms his own expectations on grounds, rather than one whose states are pushed into him from without. To produce trust by the inferential cause is to treat the one in whom the trust arises as a mind that judges, not as a mechanism that registers. And the mechanical cause, applied to a mind, does not improve upon this; it degrades it, for to compel an expectation rather than to ground it is to treat the mind, in just that respect, as if it were not a mind at all.

\subsection{Trust addresses the other as free}
\label{p13:sub:free-asfree}

The same distinction, looked at from the side of the one who trusts rather than the one who is trusted, is the distinction between two stances one may take toward another person, and it is here that Kant's worry from the diagnosis (\xref{p13:sub:diag-humekant}) receives its positive answer. One may take toward another the stance of prediction-and-management: treat his future conduct as an outcome to be forecast from evidence and secured by incentive, arrange the consequences so that the conduct one wants is the conduct it is in his interest to produce, and rely upon him exactly as one relies upon any well-understood part of the causal order. This is the stance of Holmes's observer toward the bad man (\xref{p13:sub:jur-holmes}), and it is, within its proper domain, entirely legitimate; one had better be able to predict and to manage. But it is a stance that takes its object as an object\,---\,as a system whose behaviour is determined by conditions one may arrange\,---\,and to take \emph{only} this stance toward a person, to model him exhaustively as such a system, is to have ceased, in Kant's terms, to treat him as an end, and so to have ceased to relate to him as a person at all.

The other stance is the stance of trust proper, and it is the stance of one free being toward another. To trust a person, in the full sense the promise calls for, is not to predict his conduct as one predicts the weather but to rely upon his freely binding himself\,---\,to address him as one who can commit himself to a future and can keep faith with that commitment because he has bound himself to it, not because the consequences have been arranged to leave him no alternative. The one who makes a genuine promise is not forecasting his own future behaviour from the inside, as one might forecast a habit; he is committing himself to it, performing the distinctively free act of placing himself under an obligation he could violate and undertaking that he will not. And the one who trusts the promise is wagering, in part, upon that freedom: relying not upon a mechanism that cannot fail but upon a free being who could fail and undertakes not to. The wager on the other's keeping of faith is, irreducibly, a wager on his freedom\,---\,and this, which the diagnosis marked as the point at which prediction-and-management changes its object and stops speaking of promising at all, is now seen to be not a softness in the account of trust but its very condition. There is no trusting a being one regards as a mere mechanism; there is only predicting it. Trust begins exactly where prediction is refused as the whole truth, and the other is met as one who is free to keep or to break faith and is relied upon to keep it.

\subsection{A word on the cognitive basis}
\label{p13:sub:free-cognitive}

That this is a metaphysical and ethical claim does not make it a claim against the natural order, and a word is owed, in anticipation of \xref{p13:sec:mechanism}, on its cognitive basis, lest the inferential cause seem to require a faculty outside nature. It does not. The forming of expectations on grounds is not an exotic operation a mind performs on rare and solemn occasions; it is the ordinary and incessant business of a brain, which is, on the account the next section develops, a predictive organ before it is anything else, continuously generating expectations of what is to come and revising them against what arrives. Trust, on that account, is not a special faculty superadded to cognition but a particular case of cognition's standing activity: a high-confidence forward expectation about the conduct of another, formed on grounds and held open to revision. The inferential cause has, that is, a perfectly natural implementation; what distinguishes it from the mechanical cause is not that it floats free of the causal order but that, within that order, it operates by furnishing grounds to a system that forms its own expectations, rather than by compelling a system that merely registers what is done to it. The metaphysics of freedom and the neuroscience of prediction are not at odds; the one is the significance of what the other describes. The section leaves the mechanism to its own chapter and returns to the consequence.

\subsection{Why the non-coercive promise is the proper case}
\label{p13:sub:free-proper}

The consequence is the one the Prelude felt before it could say it, and it can now be said. The promise that carries no force is not the defective limiting case of promising, to be apologised for against the enforceable contract as the standard; it is the proper case, of which the enforceable contract is the derivative. For if trust addresses the other as free, and if to produce trust is to ground a free mind's expectation rather than to compel a mechanism's response, then the promise that relies on no sword is the promise that does this in its purity: it asks to be believed on no ground but the promisor's free commitment and the promisee's free wager upon it, and so it is the form in which the inferential cause operates unmixed. The enforceable promise, by contrast, mixes into this pure case an element of the mechanical: it offers, alongside the free commitment, a threatened consequence, and to exactly that extent it addresses the promisee not as a free mind forming an expectation on the ground of trust but as a calculator forming a prediction on the ground of incentive. The sword does not strengthen the promise as a promise; it supplements the promise with something that is not a promise at all, a managed consequence, and the supplement does its work, where it works, by treating the parties in the very manner that trust refuses.

This is not yet to say\,---\,and the paper will be at pains, later, not to say\,---\,that the sword should be dispensed with, or that the enforceable contract is in every respect inferior to the bare vow; the protection of the vulnerable, as \xref{p13:sec:judicial} will insist, has a continuing and indispensable need of force. It is to say something narrower and deeper: that in respect of what makes a promise a promise\,---\,the addressing of a free being as free, the grounding rather than the compelling of an expectation\,---\,the non-coercive promise is the realised case and the coercive one the diluted. The want of force that the Prelude felt as the condition of the vow's being what it was is here given its reason. A bond one could be forced to keep is, just insofar as one could be forced, not yet fully a bond between free beings; and the vow that asks for no sword is the truer not despite carrying no force but because, carrying none, it leaves the whole of the work to freedom\,---\,to the promisor's free binding of himself and the promisee's free trust in it. How that work is in fact done\,---\,by what mechanism a free mind, addressed as free and given the right grounds, comes to expect that the promise will be kept\,---\,is the question the paper has been clearing the ground to ask, and it is the question to which it now turns.
